\chapter{Introduction}
\label{chap:introduction}

\section{Background}
\label{section:background}

% <TIP: 
% A background chapter in a book serves the purpose of providing essential context, information, or history that is relevant
% to the overall understanding of the book's subject matter. This chapter typically appears early in the book and aims to set the
% stage for the reader by offering background details that are crucial for comprehending the main narrative.
% />
The resolution of civil disputes in Thailand relies heavily on the interpretation and application of the Thai Civil and Commercial Code, alongside historical rulings from the Supreme Court. Traditionally, evaluating the merits of a civil case, formulating a legal strategy, and preparing for trial requires extensive manual effort. Legal professionals must synthesize complex factual matrices, identify applicable statutory provisions, and anticipate opposing arguments. Similarly, for legal education, practical experience is largely constrained to static case studies or resource-intensive moot court exercises. This leaves a critical gap in accessible, real-time procedural practice, creating a barrier to both dynamic legal training and cost-effective preliminary case evaluation.
\section{Problem Statement}
\label{section:problem-statement}

% <TIP: 
% A problem statement refers to a clear central issue,
% challenge, or question that the project aims to address or explore. It is a
% declaration that highlights the specific problem the author intends to examine,
% discuss, or solve throughout the course of the project.
% />
The current ecosystem for Thai civil law education, case evaluation, and trial preparation faces several critical inefficiencies:
\begin{enumerate}
    \item \textbf{The ``Experience Gap'' in Legal Education:} Law students spend years studying static texts, such as the Thai Civil and Commercial Code (TCCC) and historical Supreme Court rulings. However, they lack accessible, dynamic environments to practice adversarial legal reasoning. Organizing traditional moot courts is highly resource-intensive, leaving students without a 24/7, risk-free sandbox to stress-test legal arguments against responsive opponents.
    \item \textbf{Resource-Intensive Preliminary Case Evaluation:} For law firms, evaluating the viability of a new civil case drains valuable billable hours. Legal professionals must synthesize complex factual matrices, identify applicable statutory provisions, and anticipate adversarial counter-arguments under tight timelines. This makes case intake a high-risk, time-consuming bottleneck that heavily strains limited staff capacity.
    \item \textbf{High Financial Barriers to Justice:} The cost of retaining human legal counsel merely for a preliminary case assessment is substantial. This financial threshold often proves prohibitive for ordinary citizens or small enterprises involved in minor to moderate civil disputes, causing them to abandon potentially valid legal claims due to the initial cost of evaluation.
    \item \textbf{Lack of an Automated Adversarial Testing Ground:} Across all stakeholders, there is no automated, low-cost simulation environment. Practitioners and students lack a tool to interactively simulate opposing arguments and observe how dynamically introduced evidence or statutory interpretations might shift the probability of a trial's outcome prior to formal litigation.
\end{enumerate}

\section{Solution Overview}
\label{section:solution-overview}

% <TIP: 
% A software solution overview provides a high-level and
% concise description of a software product or system. It serves as an
% introduction to the software, offering a glimpse into its key features,
% functionalities, and the problems it aims to address. This overview is often
% presented in documentation, marketing materials, or other communication
% channels to give stakeholders, potential users, or decision-makers a quick
% understanding of what the software does and why it is valuable.
% />
\textbf{ChatreeDeka} is an AI-driven software platform designed to simulate Thai civil courtroom proceedings and evaluate the viability of civil disputes. By leveraging Multi-Agent Reinforcement Learning (MARL), the system provides a highly interactive, conversational environment that strictly aligns with the Thai Civil and Commercial Code. 

The primary objective of the software is to democratize early-stage legal assessment and provide a risk-free training ground for legal professionals and students. It directly addresses the traditional barriers of time, high financial costs, and statutory complexity by allowing users to stress-test legal arguments, introduce dynamic evidence, and receive objective, data-driven predictions on trial outcomes before committing to formal litigation.

\subsection{Features}
\label{subsection:features}

\begin{enumerate}[leftmargin=80pt]
    \item \textbf{Dynamic Trial Simulation:} Simulates the courtroom proceeding in a real-time chat interface where users and AI take turns presenting arguments, examining witnesses, and defending with counter-arguments and cross-examination questions. Users may also raise legal objections during the trial (if acting as counsel) or rule on them (if acting as the judge).
    \item \textbf{Case Setup \& Role Configuration:} Provides a conversational, jargon-free interface for non-legal users to easily input their real-world civil dispute facts, and allows users to assign manual or AI control to the Judge, Prosecutor, and Defender roles before starting a simulation.
    \item \textbf{Real-Time Evidence Injection:} Enables users to introduce new facts or evidence midway through the trial, prompting the AI agents to dynamically adjust their legal strategies.
    \item \textbf{Post-Trial Analysis \& Recommendations:} Analyzes the complete trial transcript, evidence, AI-generated content, and the adjudicating panel's profiles to calculate and display a final win probability percentage for each party. It also automatically detects low-value or high-risk cases based on the simulation outcome and suggests appropriate mediation or out-of-court settlement options.
    % \item \textbf{Dynamic Chat-Based Trial:} Simulates the courtroom proceeding in a real-time chat interface where users and AI take turns presenting arguments, examining witnesses, and defending with counter-arguments and cross-examination questions.
    % \item \textbf{Post-Trial Case Predictor:} Analyzes the complete trial transcript, evidence, AI-generated content, and the adjudicating panel's profiles to calculate and display a final win probability percentage for each party.
    % \item \textbf{Real-Time Evidence Injection:} Enables users to introduce new facts or evidence midway through the trial, prompting the AI agents to dynamically adjust their legal strategies.
    % \item \textbf{Hybrid Role Configuration:} Allows users to assign manual or AI control to the Judge, Prosecutor, and Defender roles before starting a simulation.
    % \item \textbf{Procedural Objection System:} Allows users to raise legal objections during the trial (if acting as counsel) or rule on them (if acting as the judge).
    % % \item \textbf{Precedent Matching Engine:} Extracts key legal themes from the completed simulation and suggests relevant past Thai Supreme Court rulings for further study or strategy building.
    % \item \textbf{Plain-Language Case Setup:} Provides a conversational, jargon-free interface for non-legal users to easily input their real-world civil dispute facts.
    % \item \textbf{ADR Routing Recommendation:} Automatically detects low-value or high-risk cases based on the simulation outcome and suggests appropriate mediation or out-of-court settlement options.
\end{enumerate}

\section{Target User}
\label{section:target-user}

% <TIP: 
% The target user in a software project refers to the specific
% group or demographic of individuals for whom the software is designed and
% developed. Identifying the target user is a crucial step in the software
% development process as it helps the development team tailor the software to
% meet the needs, preferences, and requirements of that particular user group.
% Understanding the characteristics, behaviors, and expectations of the target
% users is essential for creating a user-friendly and effective software solution.

% Here are some key aspects related to defining the target user in a software project:

% Demographics: This includes factors such as age, gender,
% occupation, education level, and other demographic characteristics. Different
% age groups or professional backgrounds may have distinct preferences and
% requirements when it comes to software usability.

% Skill Level: Consideration of the users' technical proficiency and
% familiarity with similar software or technology. The level of technical expertise
% can influence the complexity of the user interface, the need for tutorials or
% documentation, and other user support features.

% Industry or Domain: For software solutions designed for specific
% industries or domains, understanding the unique challenges, workflows, and
% terminology within that industry is crucial. Tailoring the software to meet
% industry-specific needs is often necessary.
% />

\textbf{ChatreeDeka} platform is designed to serve a diverse user base, ranging from highly trained legal professionals to the general public. Understanding the distinct characteristics of these groups is essential for delivering an intuitive, accessible, and effective simulation experience.

\subsection{Law Students}
\begin{itemize}
    \item \textbf{Demographics:} University students and recent graduates, typically aged 18--25, enrolled in undergraduate or graduate law programs.
    \item \textbf{Skill Level:} High technological literacy, comfortable with chat-based interfaces and dynamic software. Their legal proficiency is beginner to intermediate; they understand legal theory but lack practical courtroom experience. 
    \item \textbf{Industry/Domain:} Legal Education. They require an interactive, forgiving environment to practice case-building, oral argumentation, and receive automated, educational feedback on their procedural objections and legal documents.
\end{itemize}

\subsection{Lawyers and Legal Practitioners}
\begin{itemize}
    \item \textbf{Demographics:} Working professionals, typically aged 25--60, possessing formal law degrees and active legal licenses.
    \item \textbf{Skill Level:} Highly proficient in Thai Civil Law and court procedures. They have moderate to high technological literacy. They expect advanced functionality, such as stress-testing legal arguments against AI opponents and interpreting detailed win-probability metrics.
    \item \textbf{Industry/Domain:} Thai Legal Services and Litigation. Their primary goal is risk assessment, case strategy optimization, and matching facts with Supreme Court precedents before officially representing a client in court.
\end{itemize}

\subsection{Judiciary Trainees and Court Staff}
\begin{itemize}
    \item \textbf{Demographics:} Legal professionals who have passed or are preparing for the rigorous Thai judicial examinations, typically aged 25--45.
    \item \textbf{Skill Level:} Advanced legal knowledge with a hyper-focus on civil procedure, evidence admissibility, and trial management. They require high-fidelity simulations to practice acting as a judge while the system simulates the opposing lawyer roles.
    \item \textbf{Industry/Domain:} The Thai Justice System and Court Administration. They are focused on maintaining courtroom decorum, ruling on objections accurately, and ensuring trial outcomes align strictly with statutory law.
\end{itemize}

\subsection{Litigants / General Public}
\begin{itemize}
    \item \textbf{Demographics:} A broad, highly diverse demographic encompassing any adult citizen facing a civil dispute (e.g., contract breach, property dispute). Age, education, and occupation vary widely.
    \item \textbf{Skill Level:} Little to no legal proficiency. Tech literacy varies greatly. They are easily overwhelmed by legal jargon and complex data-entry forms. They require a plain-language, conversational interface to input their case facts.
    \item \textbf{Industry/Domain:} Personal Dispute Resolution. Their primary goal is to understand their basic legal standing, view simplified case outcome predictions, and explore Alternative Dispute Resolution (ADR) options before incurring actual legal fees.
\end{itemize}

\section{Terminology}
\label{section:terminology}

% <TIP: 
% Terminology refers to the specific language, jargon, or
% specialized vocabulary used to describe concepts, ideas, or subjects within a
% particular field or domain. The use of terminology is often essential for clarity
% and precision, especially in books that cover technical, scientific, academic, or
% specialized topics.
% />


\begin{description}

\item[TCCC (Thai Civil and Commercial Code)] The primary body of laws used to govern the simulated disputes.

\item[Matra (Section)] The specific, numbered legal rules within the Thai legal code.

\item[Deka (Dika)] A final ruling by the Supreme Court of Thailand, used in the system as the ``Ground Truth'' to evaluate AI performance.

\item[Plaintiff] The AI agent initiating the lawsuit, seeking remedy or damages.

\item[Defendant] The AI agent responding to the lawsuit, generating counter-arguments and defenses.

\item[Judge] The neutral arbiter (often modeled stochastically) that rules on motions, evaluates evidence, and determines case outcomes.

\item[Adversarial Multi-Agent] A system setup where two AI models are programmed with conflicting goals (e.g., prosecution vs.\ defense) to simulate a courtroom battle.

\item[RAG (Retrieval-Augmented Generation)] The architecture that allows AI agents to search and retrieve actual Thai laws from a vector database before formulating their arguments.

\item[MARL (Multi-Agent Reinforcement Learning)] The machine learning framework where agents (litigants, judge) learn optimal strategies by interacting within the shared simulation environment.

\item[Self-Play Training] A training method where agents improve and discover legal strategies by continuously competing against previous versions of themselves.

\item[Action Space] The specific set of legal actions available to the AI agents at any given time (e.g., filing a motion, requesting discovery, raising an objection).

\item[Policy ($\pi$)] The AI's learned strategy---essentially how it maps the current state of the trial to its next chosen legal action to maximize its chance of winning.

\item[Reward Function ($R$)] The feedback signal used to train the AI, usually rewarding victories (+1) while penalizing losses or procedural errors.

\item[Stochastic Judge Model] A probabilistic model used for the Judge agent, designed to make decisions based on real-world empirical grant rates rather than hardcoded logic.

\item[Hallucination] When the AI incorrectly fabricates a Thai law section, fact, or precedent that does not actually exist in the provided legal database.

\end{description}